\begin{textsample}{POS Dim 2 – unprofiled\_gpt – Score 8.00 – unprofiled\_gpt/t000318\_gpt.txt}  \label{ex:f2_pos_023}
I keep coming back to the \textbf{sense} that what we ’re living through now really started in that Thatcher–Reagan era, when this particular brand of neoliberal capitalism got unleashed.
It feels like they lit the fuse on something that ’s just rolled on and on, reshaping everything around markets and profit, and it ’s done a \textbf{lot} of damage to the world.
Sometimes I find myself \textbf{almost} blaming them personally for \textbf{setting} us on this course.
But then I \textbf{also} wonder if there ’s a strange kind of necessity in it, you know?
That maybe things have to swing so \textbf{far} towards selfishness and greed before people collectively say, hang on, this is intolerable.
It ’s like you need that extreme to provoke a proper global reaction.
And now you ’ve got people like Corbyn back in the UK and that fella in the States, Bill Sanders, standing up and saying, look, we all have to take responsibility, we have to fight corruption, we can't just leave it to the same old political class.
It makes me think of the late eighties, when communism suddenly crumbled.
For years it looked immovable, and then \textbf{almost} overnight it was gone.
I can't help feeling we might be at a similar kind of turning point with this current system.
Maybe it has to get as bad as it possibly can before it just collapses under its \textbf{own} weight and something new emerges.
There ’s something quite momentous about that \textbf{idea}, that we might be watching the end of a whole way of organising the world, even if we don't yet know what comes next.

% matched lemmas: almost, also, far, idea, lot, own, sense, set
\end{textsample}
